\section{Robust Partitions via Covering}
\label{sec:partitions-covering}

In this section, we pursue two goals. First, we introduce a formulation of $(\mu,\epsilon)$-robustness using partitions. We then show that $(\mu,\epsilon)$ is equivalent to covering, mirroring equivalence of $\epsilon$-robustness to separability.

\subsection{Robust Partitions}

To avoid ambiguities, we introduce some well-known definitions from the literature that are relevant to our developments.
% lax begin Lax544225.Partitions
\begin{definition}[Partition]
\label{def:partition}
A collection of nonempty sets $\{P_1,\dots,P_p\}$ is a partition of a set $S$ iff
\begin{enumerate}
    \item $\bigcup_{i \in [p]} P_i = S$ and
    \item $P_i \cap P_j = \emptyset$ for all $i \neq j$.
\end{enumerate} 
The partition $\{S\}$ on the set $S$ is called \emph{trivial}.
\end{definition}
% lax end

% lax begin Lax544225.Partitions
\begin{definition}[Classifier Induced Partition]
\label{def:induced-partition}
    A classifier $c: S \rightarrow [n]$ induces a partition $P^c := \{S^c_1, S^c_2, \dots S^c_n\}$ with $S^c_i := \{x \in S\,|\,c(x) = i\}$.
\end{definition}
% lax end

Robustness is often described as invariance under $\epsilon$-sized perturbations. This perspective allows us to reformulate robustness, for which we define the close neighborhood around an input as an $\epsilon$-ball:
% lax begin Lax544225.Balls
\begin{definition}[$\epsilon$-Ball]
\label{def:epsilon-ball}
Let $(S, d)$ be a semimetric space. Then $B(x^*, \epsilon) := \{x \in S\,|\,d(x,x^*) \leq \epsilon\}$ is an \emph{$\epsilon$-ball} around $x^*$.
\end{definition}
% lax end

Despite the name, $B(x^*, \epsilon)$ may not resemble a geometric ball. In fairness notions, for example, the distance function may capture similarity in qualification of individuals rather than geometric proximity of their attribute vectors.
% lax begin Lax544225.PartitionRobustness
\begin{definition}[$(\mu,\epsilon)$-Robustness of a Partition]
\label{def:mueps-robustness-partition}
Let $(S,d)$ be a semimetric space. Then a partition $P := \{S_1,\dots,S_p\}$ of $S$
is \emph{globally robust} for $\mu \geq 0, \epsilon > 0$ if and only if:
$$\forall x \in S\,\exists x' \in B(x,\mu) \,[\bigvee_{i \in [p]} (\{x\} \cup B(x', \epsilon) \subseteq S_i)].$$
\end{definition}
% lax end

This subtle reformulation completely shifts $(\mu,\epsilon)$-robustness to a property of the induced partition. 
By working with partitions, our framework does not require the number of classes to be specified a priori, unlike some existing approaches~\cite{yang2024geometric}. To streamline subsequent results, we will use the term \emph{$(\mu,\epsilon)$-robust partition} for partitions that satisfy Definition~\ref{def:mueps-robustness-partition}. If parameters are clear from the context, we will omit $(\mu,\epsilon)$.

\subsection{From Separability to Covering}

A folklore fact about global robustness is its equivalence to \emph{$\epsilon$-separability}~\cite{yang2020closer,mangal2023certifying} and we can to extend this property to our new robustness notion, albeit in a more complicated form. 
% lax begin Lax544225.CoveringCharacterization
\begin{theorem}
\label{thm:separability}
    Let $(S,d)$ a semimetric space, $P := \{S_1, \dots, S_n\}$ a partition of $S$ and $\mu \geq 0, \epsilon > 0$. Then the following are equivalent:
% lax end
    \begin{enumerate}
        \item $P$ is $(\mu,\epsilon)$-robust.
% lax begin Lax544225.Covering
        \item for all $S_i \in P$ and all $x \in S_i$ holds that $B(x, \mu) \not\subseteq \bigcup_{x^* \in S \setminus S_i} B(x^*, \epsilon)$.
% lax end

    \end{enumerate}
\end{theorem}
\begin{proof}
    We prove both directions separately.

    $(1.) \Rightarrow (2.)$: 
    For contradiction, assume that (1.) holds but (2.) does not. In this case, there exists $S_i \in P$ with $x \in S_i$ such that $B(x,\mu) \subseteq \bigcup_{x^* \in S \setminus S_i} B(x^*, \epsilon)$. From Defintion~\ref{def:mueps-robustness-partition}, there exists $x' \in B(x,\mu)$ with $B(x',\epsilon) \subseteq S_i$. From symmetry of distance (Definition~\ref{def:distance}), it follows that $x' \notin \bigcup_{x^* \in S \setminus S_i} B(x^*, \epsilon)$, leading to a contradiction with (1.).

    $(2.) \Rightarrow (1.)$:
    Similarly, assume that (2.) holds but not (1.). By Defintion~\ref{def:mueps-robustness-partition}, there exists $x \in S$ such that for all $x' \in B(x,\mu)$ have $B(x', \epsilon) \not\subseteq S_i$ for all $S_i \in P$. Pick a specific $S_i' \in P$ such that $x \in S_i'$. By symmetry of distance, each $x' \in B(x,\mu)$ is also contained in some $B(x^*,\epsilon)$ with $x^* \notin S_i'$. This implies $B(x, \mu) \subseteq \bigcup_{x^* \in S \setminus S_i'} B(x^*, \epsilon)$, raising a contradiction to (2.).
    \qed
\end{proof}

\begin{figure}[t]
  \centering
  \definecolor{muBlue}{HTML}{1F4E79}
  \definecolor{epsGreen}{HTML}{2E8B57}
  \definecolor{epsOrange}{HTML}{D9730D}
  \begin{tikzpicture}[scale=1.5, every node/.style={transform shape},
      classLine/.style = {black!75, line width=0.9pt},
      ballS2/.style    = {draw=epsGreen,  thick},
      ballS3/.style    = {draw=epsOrange, thick},
      dotS1/.style     = {circle, fill=muBlue,    inner sep=1.5pt},
      dotS2/.style     = {circle, fill=epsGreen,  inner sep=1.5pt},
      dotS3/.style     = {circle, fill=epsOrange, inner sep=1.5pt},
      radArrow/.style  = {-{Stealth[length=2.2mm]}, thin, gray!70},
    ]
    % ----------------------------------------------------------------
    % Geometry (numerically verified):
    %   decision boundary (vertical, wiggly):
    %     x(y) = 0.55 - 0.85*cos(3y)   [3 rad = 171.9 degrees]
    %   S_1 = right of the curve, S_2/S_3 = left of it (split at y=-0.1)
    %   x   = (0, 0) in S_1,  mu  = 1.1
    %   eps = 1.0   (mu = 1.1*eps > eps)
    %   eps-balls (radius 1.0), centers in S_2 / S_3:
    %     x_1 = ( 0.28,  0.48)  in S_2
    %     x_2 = (-0.55,  0.00)  in S_2
    %     x_3 = ( 0.28, -0.48)  in S_3
    %   The three eps-balls cover the mu-ball of x completely
    %   (max distance of any mu-ball point to the nearest center: 0.953);
    %   neither S_2 alone nor S_3 alone covers it.
    %   Note: three eps-disks can cover a mu-disk of radius at most
    %   2*eps/sqrt(3) ~ 1.155*eps, so mu = 1.1*eps is near the
    %   largest ratio possible with three covering balls.
    %   Canvas is cropped tight: x in [-1.75,1.5], y in [-1.7,1.7].
    % ----------------------------------------------------------------
    \coordinate (x)  at (0,0);
    \coordinate (x1) at (0.28,0.48);
    \coordinate (x2) at (-0.55,0);
    \coordinate (x3) at (0.28,-0.48);

    % ---- landscape: the three class regions ---------------------------
    \fill[blue!8]                                            % S_1 (right)
      plot[domain=1.7:-1.7, samples=250, smooth]
        ({0.55-0.85*cos(171.9*\x)},{\x})
      -- (1.5,-1.7) -- (1.5,1.7) -- cycle;
    \fill[green!12]                                          % S_2 (upper left)
      (-1.75,1.7)
      -- plot[domain=1.7:-0.1, samples=150, smooth]
           ({0.55-0.85*cos(171.9*\x)},{\x})
      -- (-1.75,-0.1) -- cycle;
    \fill[orange!10]                                         % S_3 (lower left)
      (-1.75,-0.1)
      -- plot[domain=-0.1:-1.7, samples=150, smooth]
           ({0.55-0.85*cos(171.9*\x)},{\x})
      -- (-1.75,-1.7) -- cycle;

    % ---- coverage blankets: union of the eps-balls per class ----------
    % (transparency group per class => each class gets ONE uniform tint)
    \begin{scope}[transparency group, opacity=0.2]
      \fill[epsGreen] (x1) circle (1.0);
      \fill[epsGreen] (x2) circle (1.0);
    \end{scope}
    \begin{scope}[transparency group, opacity=0.2]
      \fill[epsOrange] (x3) circle (1.0);
    \end{scope}

    % ---- class boundaries ----------------------------------------------
    \draw[classLine]
      plot[domain=-1.7:1.7, samples=250, smooth]
        ({0.55-0.85*cos(171.9*\x)},{\x});
    \draw[classLine] (-1.75,-0.1) -- (-0.26,-0.1);

    % ---- highlighted epsilon-balls ----------------------------------------
    \draw[ballS2] (x1) circle (1.0);
    \draw[ballS2] (x2) circle (1.0);
    \draw[ballS3] (x3) circle (1.0);

    % ---- mu-ball around x (mu > eps) --------------------------------------
    \draw[muBlue, very thick] (x) circle (1.1);

    % ---- points --------------------------------------------------------------
    \node[dotS1] at (x)  {};
    \node[dotS2] at (x1) {};
    \node[dotS2] at (x2) {};
    \node[dotS3] at (x3) {};

    % ---- radius annotations -----------------------------------------------------
    \draw[radArrow] (x)  -- ($(x)+(250:1.1)$)  node[midway, left]        {$\mu$};
    \draw[radArrow] (x1) -- ($(x1)+(60:1.0)$)  node[midway, above right] {$\epsilon$};
    \draw[radArrow] (x2) -- ($(x2)+(150:1.0)$) node[midway, above left] {$\epsilon$};
    \draw[radArrow] (x3) -- ($(x3)+(300:1.0)$) node[midway, below right] {$\epsilon$};

    % ---- point labels ----------------------------------------------------------------
    \node[below right=0pt]  at (x)  {$x$};
    \node[above right=1pt]  at (x1) {$x_1$};
    \node[left=2pt]         at (x2) {$x_2$};
    \node[below right=1pt]  at (x3) {$x_3$};

    % ---- region labels --------------------------------------------------------------------
    \node[font=\small] at (-1.5, 1.45) {$S_2$};  % upper left
    \node[font=\small] at (-1.5,-1.3)  {$S_3$};  % lower left
    \node[font=\small] at ( 1.25, 1.45) {$S_1$}; % right
  \end{tikzpicture}
  \caption{Illustrated violation of $(\mu,\epsilon)$-robustness though covering condition (2.) of Theorem~\ref{thm:separability}: the $\mu$-ball of $x \in S_1$ is fully covered by three $\epsilon$-balls of $x_1, x_2 \in S_2$ and $x_3 \in S_3$.}
  \label{fig:covering}
\end{figure}

The condition (2.) in Theorem~\ref{thm:separability} can be stated as: no $\mu$-neighbourhood is covered by $\epsilon$-neighbourhoods of points with a different label. An example when it is not the case is illustated in Figure~\ref{fig:covering}. As $(\mu, \epsilon)$-robustness is a generalization of $\epsilon$-robustness, our proof is also a formalization of separability for $\epsilon$-robustness. For $\mu = 0$, the condition (2.) reduces to each point (0-neighbourhood) being $\epsilon$-far from any point with a different label. Separability is established here in a general, non-probabilistic setting without distributional assumptions. Unlike prior work, we show that separability is a \emph{necessary} condition for global robustness. 

The characterization of robustness from Theorem~\ref{thm:separability} is central to future results and we will use it as a primary criterion of partition robustness.

